% ============================================================
\chapter{Modul 10: Fehlerbehandlung}
% ============================================================

Sensoren liefern falsche Werte. Verbindungen brechen ab. Dateien fehlen. Ein Roboter, der bei jedem Fehler abstürzt, ist nutzlos. Fehlerbehandlung ist kein optionales Extra -- sie ist Sicherheitspflicht.

\section{Lektion 29: Exceptions}

\begin{lstlisting}
import serial   # Serielle Kommunikation

# OHNE Fehlerbehandlung -- gefaehrlich!
# port = serial.Serial("/dev/ttyUSB0", 9600)  # Absturz wenn Port fehlt

# MIT Fehlerbehandlung -- robust
try:
    port = serial.Serial("/dev/ttyUSB0", 9600, timeout=1)
    daten = port.readline().decode("utf-8").strip()
    wert  = float(daten)
    print(f"Sensorwert: {wert}")

except serial.SerialException as e:
    print(f"FEHLER: Port nicht verfuegbar -- {e}")
    wert = None

except ValueError as e:
    print(f"FEHLER: Ungueltige Daten '{daten}' -- {e}")
    wert = None

except Exception as e:
    print(f"UNBEKANNTER FEHLER: {e}")
    wert = None

finally:
    print("Verbindungsversuch abgeschlossen.")
    # finally laeuft IMMER -- ideal fuer Aufraeumarbeiten
\end{lstlisting}

\begin{merksatz}
\texttt{try} -- ``versuche das''. \texttt{except TypError} -- ``falls dieser Fehler auftritt''. \texttt{finally} -- ``das passiert immer, egal was''. Mehrere \texttt{except}-Blöcke fangen verschiedene Fehlertypen ab. Vom spezifischsten zum allgemeinsten.
\end{merksatz}

\section{Lektion 30: Robuste Programme}

\begin{lstlisting}
class RobusterSensor:
    """Sensor mit eingebautem Fehler-Fallback."""

    def __init__(self, port, max_versuche=3):
        self.port          = port
        self.max_versuche  = max_versuche
        self.fehleranzahl  = 0
        self.letzter_wert  = None

    def lese(self):
        for versuch in range(self.max_versuche):
            try:
                # Simulierter Sensorlesefehler
                import random
                if random.random() < 0.3:   # 30% Fehlerrate
                    raise ValueError("Pruefsummenfehler")
                wert = round(random.uniform(0.1, 2.0), 2)
                self.letzter_wert = wert
                self.fehleranzahl = 0
                return wert

            except ValueError as e:
                self.fehleranzahl += 1
                print(f"  Versuch {versuch+1}/{self.max_versuche}: {e}")

        print(f"WARNUNG: Sensor {self.port} nach {self.max_versuche} Versuchen nicht erreichbar.")
        return self.letzter_wert   # Letzten bekannten Wert zurueckgeben

    def ist_ok(self):
        return self.fehleranzahl < self.max_versuche


sensor = RobusterSensor("ultraschall")
for i in range(10):
    wert = sensor.lese()
    if wert is not None:
        print(f"Abstand: {wert} m")
\end{lstlisting}

\begin{praxis}
Echte Roboter nutzen dieses Muster überall: drei Versuche, dann Fallback. In sicherheitskritischen Systemen gibt es zusätzlich Watchdog-Timer: Wenn kein gültiger Sensorwert innerhalb von X Millisekunden -- Notabschaltung.
\end{praxis}

\begin{uebung}[title=Projektaufgabe Modul 10: Robuster Datenleser]
Schreibe eine Funktion \texttt{lade\_config\_sicher(pfad)}, die eine JSON-Konfigurationsdatei lädt. Behandle: Datei nicht gefunden (\texttt{FileNotFoundError}), ungültiges JSON (\texttt{json.JSONDecodeError}), fehlende Pflichtparameter. Bei Fehler: Standardkonfiguration zurückgeben und Fehler loggen.
\end{uebung}
